\chapter{Communication Complexity — Public Coin Composition}
%%% generated by the blueprint pipeline from TCSlib.CommunicationComplexity.NewmanTheorem.PublicCoinComposition
\section{Overview}

This module defines composition operations on public-coin finite-message communication
protocols: output mapping, sequential binding, randomness reindexing, pairwise and
$k$-fold products, fresh-randomness binding, and mixed public-coin/deterministic products.
For each combinator it establishes how running the composed protocol reduces to running the
components, and how the communication complexity of the composed protocol relates to those
of its parts.

%%%%%%%%%%%%%%%%%%%%%%%%%%%%%%%%%%%%%%%%%%%%%%%%%%%%%%%%%%%%%%%%%%%%%%%%%%%%%%%
\section{Declarations}

\begin{definition}[Map over protocol output]
\lean{CommunicationComplexity.PublicCoin.FiniteMessage.Protocol.map}
\label{CommunicationComplexity.PublicCoin.FiniteMessage.Protocol.map}
\leanok
\uses{CommunicationComplexity.Deterministic.FiniteMessage.Protocol.map, CommunicationComplexity.PublicCoin.FiniteMessage.Protocol}
Given a function $g : \alpha \to \beta$ and a public-coin protocol
$p : \mathrm{Protocol}\;\Omega\;X\;Y\;\alpha$, the protocol $p.\mathrm{map}\;g$
is a $\mathrm{Protocol}\;\Omega\;X\;Y\;\beta$ that applies $g$ to the output of $p$
after every run, without changing the communication structure.
\end{definition}

\begin{theorem}[Mapping the output commutes with randomized execution]
\lean{CommunicationComplexity.PublicCoin.FiniteMessage.Protocol.map_rrun}
\label{CommunicationComplexity.PublicCoin.FiniteMessage.Protocol.map_rrun}
\leanok
\uses{CommunicationComplexity.Deterministic.FiniteMessage.Protocol.map_run, CommunicationComplexity.PublicCoin.FiniteMessage.Protocol, CommunicationComplexity.PublicCoin.FiniteMessage.Protocol.map, CommunicationComplexity.PublicCoin.FiniteMessage.Protocol.rrun}
\difficulty{1}
Let $p$ be a public-coin finite-message protocol over shared-randomness space $\Omega$,
private input sets $X$ and $Y$, and output type $\alpha$, and let
$g\colon\alpha\to\beta$ be any function. Write $\mathrm{map}(g,p)$ for the protocol with
the same communication structure as $p$ that applies $g$ to whatever value $p$ outputs,
and write $\mathrm{rrun}(q,x,y,\omega)$ for the result of executing a protocol $q$ on
private inputs $x$ and $y$ with shared randomness $\omega$. Then for all $x\in X$, $y\in
Y$, and $\omega\in\Omega$,
\[
\mathrm{rrun}\bigl(\mathrm{map}(g,p),\,x,\,y,\,\omega\bigr)=g\bigl(\mathrm{rrun}(p,\,x,\,y,\,\omega)\bigr).
\]
\end{theorem}

\begin{theorem}[Map over the output preserves complexity]
\lean{CommunicationComplexity.PublicCoin.FiniteMessage.Protocol.map_complexity}
\label{CommunicationComplexity.PublicCoin.FiniteMessage.Protocol.map_complexity}
\leanok
\uses{CommunicationComplexity.Deterministic.FiniteMessage.Protocol.complexity, CommunicationComplexity.Deterministic.FiniteMessage.Protocol.map_complexity, CommunicationComplexity.PublicCoin.FiniteMessage.Protocol, CommunicationComplexity.PublicCoin.FiniteMessage.Protocol.map}
\difficulty{1}
Let $p$ be a public-coin finite-message protocol over a shared-randomness space
$\Omega$, private input sets $X$ and $Y$, and output type $\alpha$, and let $g : \alpha
\to \beta$ be any function. Then the protocol $p.\mathrm{map}\,g$, which runs $p$ and
applies $g$ to its output, has the same communication complexity as $p$:
\[
\mathrm{complexity}(p.\mathrm{map}\,g) = \mathrm{complexity}(p).
\]
\end{theorem}

\begin{definition}[Bind on protocols]
\lean{CommunicationComplexity.PublicCoin.FiniteMessage.Protocol.bind}
\label{CommunicationComplexity.PublicCoin.FiniteMessage.Protocol.bind}
\leanok
\uses{CommunicationComplexity.Deterministic.FiniteMessage.Protocol.bind, CommunicationComplexity.PublicCoin.FiniteMessage.Protocol}
Given a protocol $p : \mathrm{Protocol}\;\Omega\;X\;Y\;\alpha$ and a family
$q : \alpha \to \mathrm{Protocol}\;\Omega\;X\;Y\;\beta$, the protocol $p.\mathrm{bind}\;q$
replaces each output $a$ of $p$ with the protocol $q\;a$, both sharing the same randomness
$\Omega$.
\end{definition}

\begin{theorem}[Bind commutes with randomized execution]
\lean{CommunicationComplexity.PublicCoin.FiniteMessage.Protocol.bind_rrun}
\label{CommunicationComplexity.PublicCoin.FiniteMessage.Protocol.bind_rrun}
\leanok
\uses{CommunicationComplexity.Deterministic.FiniteMessage.Protocol.bind_run, CommunicationComplexity.PublicCoin.FiniteMessage.Protocol, CommunicationComplexity.PublicCoin.FiniteMessage.Protocol.bind, CommunicationComplexity.PublicCoin.FiniteMessage.Protocol.rrun}
\difficulty{1}
Let $p$ be a public-coin finite-message protocol with shared-randomness space $\Omega$,
private input sets $X$ and $Y$, and output type $\alpha$, and let $q$ be a family of
such protocols indexed by $\alpha$ with output type $\beta$, all sharing the same
randomness space $\Omega$. Then for every private input $x \in X$, every private input
$y \in Y$, and every shared random string $\omega \in \Omega$, the randomized execution
of the sequential composition $p.\mathrm{bind}\,q$ agrees with first executing $p$ and
then executing the selected continuation:
\[
\mathrm{rrun}(p.\mathrm{bind}\,q,\, x,\, y,\, \omega)
  = \mathrm{rrun}\bigl(q(\mathrm{rrun}(p, x, y, \omega)),\, x,\, y,\, \omega\bigr).
\]
\end{theorem}

\begin{definition}[Randomness reindexing]
\lean{CommunicationComplexity.PublicCoin.FiniteMessage.Protocol.comapRandomness}
\label{CommunicationComplexity.PublicCoin.FiniteMessage.Protocol.comapRandomness}
\leanok
\uses{CommunicationComplexity.Deterministic.FiniteMessage.Protocol.comap, CommunicationComplexity.PublicCoin.FiniteMessage.Protocol}
Given $h : \Omega' \to \Omega$ and a protocol $p : \mathrm{Protocol}\;\Omega\;X\;Y\;\alpha$,
the protocol $\mathrm{comapRandomness}\;h\;p$ is a
$\mathrm{Protocol}\;\Omega'\;X\;Y\;\alpha$ that runs $p$ after pre-composing the shared
randomness with $h$.
\end{definition}

\begin{theorem}[Randomness reindexing commutes with randomized execution]
\lean{CommunicationComplexity.PublicCoin.FiniteMessage.Protocol.comapRandomness_rrun}
\label{CommunicationComplexity.PublicCoin.FiniteMessage.Protocol.comapRandomness_rrun}
\leanok
\uses{CommunicationComplexity.Deterministic.FiniteMessage.Protocol.comap_run, CommunicationComplexity.PublicCoin.FiniteMessage.Protocol, CommunicationComplexity.PublicCoin.FiniteMessage.Protocol.comapRandomness, CommunicationComplexity.PublicCoin.FiniteMessage.Protocol.rrun}
\difficulty{1}
Let $h : \Omega' \to \Omega$ be a map between shared-randomness spaces, and let $p$ be a
public-coin finite-message protocol over randomness space $\Omega$, private input sets
$X$ and $Y$, and output type $\alpha$. Form the reindexed protocol over $\Omega'$ that
pre-composes the shared randomness with $h$ before running $p$. Then for all private
inputs $x \in X$ and $y \in Y$ and all shared randomness $\omega \in \Omega'$, executing
the reindexed protocol on inputs $x, y$ with randomness $\omega$ produces the same
output as executing $p$ on inputs $x, y$ with randomness $h(\omega)$.
\end{theorem}

\begin{theorem}[Randomness reindexing preserves communication complexity]
\lean{CommunicationComplexity.PublicCoin.FiniteMessage.Protocol.comapRandomness_complexity}
\label{CommunicationComplexity.PublicCoin.FiniteMessage.Protocol.comapRandomness_complexity}
\leanok
\uses{CommunicationComplexity.Deterministic.FiniteMessage.Protocol.complexity, CommunicationComplexity.PublicCoin.FiniteMessage.Protocol, CommunicationComplexity.PublicCoin.FiniteMessage.Protocol.comapRandomness}
\difficulty{1}
Let $p$ be a public-coin finite-message protocol with shared-randomness space $\Omega$,
private inputs $X$ and $Y$, and outputs in $\alpha$, and let $h : \Omega' \to \Omega$ be
any map. Form the protocol over the new randomness space $\Omega'$ by replacing each
player's effective input $\Omega \times X$ or $\Omega \times Y$ with $\Omega' \times X$
or $\Omega' \times Y$ and precomposing every message function with $h$ on the randomness
coordinate (leaving the private inputs untouched). Then this reindexed protocol has
exactly the same worst-case communication complexity as $p$.
\end{theorem}

\begin{definition}[Product of two protocols]
\lean{CommunicationComplexity.PublicCoin.FiniteMessage.Protocol.prod}
\label{CommunicationComplexity.PublicCoin.FiniteMessage.Protocol.prod}
\leanok
\uses{CommunicationComplexity.Deterministic.FiniteMessage.Protocol.comap, CommunicationComplexity.Deterministic.FiniteMessage.Protocol.prod, CommunicationComplexity.PublicCoin.FiniteMessage.Protocol}
Given protocols $p_1 : \mathrm{Protocol}\;\Omega_1\;X\;Y\;\alpha_1$ and
$p_2 : \mathrm{Protocol}\;\Omega_2\;X\;Y\;\alpha_2$, the product protocol
$\mathrm{prod}\;p_1\;p_2$ is a $\mathrm{Protocol}\;(\Omega_1 \times \Omega_2)\;X\;Y\;(\alpha_1 \times \alpha_2)$
that runs $p_1$ and $p_2$ independently on the same inputs, using the first and second
components of the shared randomness pair respectively.
\end{definition}

\begin{theorem}[Product protocol runs component-wise]
\lean{CommunicationComplexity.PublicCoin.FiniteMessage.Protocol.prod_rrun}
\label{CommunicationComplexity.PublicCoin.FiniteMessage.Protocol.prod_rrun}
\leanok
\uses{CommunicationComplexity.Deterministic.FiniteMessage.Protocol.comap_run, CommunicationComplexity.Deterministic.FiniteMessage.Protocol.prod_run, CommunicationComplexity.PublicCoin.FiniteMessage.Protocol, CommunicationComplexity.PublicCoin.FiniteMessage.Protocol.prod, CommunicationComplexity.PublicCoin.FiniteMessage.Protocol.rrun}
\difficulty{1}
Let $p_1$ be a public-coin finite-message protocol with shared-randomness space
$\Omega_1$ and output type $\alpha_1$, and let $p_2$ be one with shared-randomness space
$\Omega_2$ and output type $\alpha_2$, both over the same private input sets $X$
(Alice's) and $Y$ (Bob's). Then for every pair of inputs $x \in X$, $y \in Y$ and every
joint random string $\omega = (\omega_1, \omega_2) \in \Omega_1 \times \Omega_2$, the
randomized execution of the product protocol $\mathrm{prod}(p_1, p_2)$ decomposes into
the executions of its factors:
\[
\mathrm{rrun}\bigl(\mathrm{prod}(p_1, p_2),\, x,\, y,\, \omega\bigr) \;=\;
\bigl(\mathrm{rrun}(p_1, x, y, \omega_1),\ \mathrm{rrun}(p_2, x, y, \omega_2)\bigr),
\]
that is, each component is obtained by running the corresponding factor on the same
inputs $x$, $y$ with the matching component of the shared randomness.
\end{theorem}

\begin{theorem}[Communication complexity of a product protocol]
\lean{CommunicationComplexity.PublicCoin.FiniteMessage.Protocol.prod_complexity}
\label{CommunicationComplexity.PublicCoin.FiniteMessage.Protocol.prod_complexity}
\leanok
\uses{CommunicationComplexity.Deterministic.FiniteMessage.Protocol.complexity, CommunicationComplexity.Deterministic.FiniteMessage.Protocol.prod_complexity, CommunicationComplexity.PublicCoin.FiniteMessage.Protocol, CommunicationComplexity.PublicCoin.FiniteMessage.Protocol.prod}
\difficulty{1}
Let $p_1$ and $p_2$ be public-coin finite-message protocols over the common private
input sets $X$ and $Y$, with shared-randomness spaces $\Omega_1$ and $\Omega_2$ and
output types $\alpha_1$ and $\alpha_2$ respectively. Then the communication complexity
of their product protocol $\mathrm{prod}\;p_1\;p_2$, which runs $p_1$ and $p_2$
independently on the same inputs, equals the sum of the individual complexities:
\[
\mathrm{complexity}(\mathrm{prod}\;p_1\;p_2) = \mathrm{complexity}(p_1) +
\mathrm{complexity}(p_2).
\]
\end{theorem}

\begin{definition}[$k$-fold product of protocols]
\lean{CommunicationComplexity.PublicCoin.FiniteMessage.Protocol.pi}
\label{CommunicationComplexity.PublicCoin.FiniteMessage.Protocol.pi}
\leanok
\uses{CommunicationComplexity.Deterministic.FiniteMessage.Protocol.comap, CommunicationComplexity.Deterministic.FiniteMessage.Protocol.pi, CommunicationComplexity.PublicCoin.FiniteMessage.Protocol}
Given a family of protocols $p : (i : \mathrm{Fin}\;k) \to \mathrm{Protocol}\;(\Omega_i)\;X\;Y\;(\alpha_i)$
with heterogeneous randomness and output types, the $\pi$-protocol
$\mathrm{pi}\;p$ is a $\mathrm{Protocol}\;(\prod_i \Omega_i)\;X\;Y\;(\prod_i \alpha_i)$
that runs each $p_i$ independently using its own component of the shared randomness tuple.
\end{definition}

\begin{theorem}[Product protocol runs component-wise]
\lean{CommunicationComplexity.PublicCoin.FiniteMessage.Protocol.pi_rrun}
\label{CommunicationComplexity.PublicCoin.FiniteMessage.Protocol.pi_rrun}
\leanok
\uses{CommunicationComplexity.Deterministic.FiniteMessage.Protocol.comap_run, CommunicationComplexity.Deterministic.FiniteMessage.Protocol.pi_run, CommunicationComplexity.PublicCoin.FiniteMessage.Protocol, CommunicationComplexity.PublicCoin.FiniteMessage.Protocol.pi, CommunicationComplexity.PublicCoin.FiniteMessage.Protocol.rrun}
\difficulty{1}
Fix $k \in \bbn$, shared private-input sets $X$ and $Y$, and, for each index $i \in
\mathrm{Fin}\,k$, a randomness space $\Omega_i$ and an output type $\alpha_i$. Let $p =
(p_i)_{i}$ be a family of public-coin finite-message protocols, where $p_i$ has
randomness space $\Omega_i$, private inputs from $X$ and $Y$, and outputs in $\alpha_i$,
and let $\mathrm{pi}\,p$ be their $k$-fold product, a public-coin protocol with
randomness space $\prod_i \Omega_i$ and output type $\prod_i \alpha_i$. Then for all
inputs $x \in X$, $y \in Y$ and every randomness tuple $\omega = (\omega_i)_i \in
\prod_i \Omega_i$, the randomized execution of the product on $(x, y, \omega)$ is the
tuple of the individual executions:
\[
\mathrm{rrun}(\mathrm{pi}\,p,\, x,\, y,\, \omega) = \bigl(i \mapsto \mathrm{rrun}(p_i,\,
x,\, y,\, \omega_i)\bigr).
\]
\end{theorem}

\begin{theorem}[Complexity of a product of public-coin protocols is additive]
\lean{CommunicationComplexity.PublicCoin.FiniteMessage.Protocol.pi_complexity}
\label{CommunicationComplexity.PublicCoin.FiniteMessage.Protocol.pi_complexity}
\leanok
\uses{CommunicationComplexity.Deterministic.FiniteMessage.Protocol.complexity, CommunicationComplexity.Deterministic.FiniteMessage.Protocol.pi_complexity, CommunicationComplexity.PublicCoin.FiniteMessage.Protocol, CommunicationComplexity.PublicCoin.FiniteMessage.Protocol.pi}
\difficulty{1}
Let $k$ be a natural number, and for each index $i \in \mathrm{Fin}\,k$ let $p_i$ be a
public-coin finite-message protocol with shared-randomness space $\Omega_i$, common
private input sets $X$ and $Y$, and output type $\alpha_i$. Then the communication
complexity of their $k$-fold product — the protocol $\mathrm{pi}\,p$ over the combined
randomness space $\prod_i \Omega_i$ that runs each $p_i$ on its own randomness component
and returns the tuple of outputs — is the sum of the complexities of the factors:
\[
\mathrm{complexity}(\mathrm{pi}\,p) \;=\; \sum_{i \in \mathrm{Fin}\,k}
\mathrm{complexity}(p_i).
\]
\end{theorem}

\begin{definition}[Bind with fresh randomness]
\lean{CommunicationComplexity.PublicCoin.FiniteMessage.Protocol.rbind}
\label{CommunicationComplexity.PublicCoin.FiniteMessage.Protocol.rbind}
\leanok
\uses{CommunicationComplexity.Deterministic.FiniteMessage.Protocol.bind, CommunicationComplexity.Deterministic.FiniteMessage.Protocol.comap, CommunicationComplexity.PublicCoin.FiniteMessage.Protocol}
Given $p : \mathrm{Protocol}\;\Omega\;X\;Y\;\alpha$ and
$q : \alpha \to \mathrm{Protocol}\;\Omega'\;X\;Y\;\beta$, the protocol
$\mathrm{rbind}\;p\;q$ is a $\mathrm{Protocol}\;(\Omega \times \Omega')\;X\;Y\;\beta$
that first runs $p$ using the $\Omega$-component of the shared randomness, then runs
$q$ on the result using the independent $\Omega'$-component.
\end{definition}

\begin{theorem}[Randomized bind unfolds into sequential execution]
\lean{CommunicationComplexity.PublicCoin.FiniteMessage.Protocol.rbind_rrun}
\label{CommunicationComplexity.PublicCoin.FiniteMessage.Protocol.rbind_rrun}
\leanok
\uses{CommunicationComplexity.Deterministic.FiniteMessage.Protocol.bind_run, CommunicationComplexity.Deterministic.FiniteMessage.Protocol.comap_run, CommunicationComplexity.PublicCoin.FiniteMessage.Protocol, CommunicationComplexity.PublicCoin.FiniteMessage.Protocol.rbind, CommunicationComplexity.PublicCoin.FiniteMessage.Protocol.rrun}
\difficulty{1}
Let $p$ be a public-coin finite-message protocol over shared-randomness space $\Omega$
with private inputs in $X$, $Y$ and outputs in $\alpha$, and let $q$ assign to each
output $a \in \alpha$ a public-coin finite-message protocol $q(a)$ over
shared-randomness space $\Omega'$ with the same private input sets and outputs in
$\beta$. Then for all private inputs $x \in X$ and $y \in Y$ and every randomness pair
$\omega = (\omega_1, \omega_2) \in \Omega \times \Omega'$, the randomized execution of
the fresh-randomness bind $\mathrm{rbind}(p, q)$ satisfies
\[
  \mathrm{rrun}\bigl(\mathrm{rbind}(p, q),\, x,\, y,\, \omega\bigr)
= \mathrm{rrun}\bigl(q(\mathrm{rrun}(p, x, y, \omega_1)),\, x,\, y,\, \omega_2\bigr).
\]
That is, running the bound protocol with randomness $\omega$ is the same as first
running $p$ on $x, y$ with the randomness $\omega_1$ to obtain an output, and then
running the protocol $q$ assigns to that output on $x, y$ with the independent
randomness $\omega_2$.
\end{theorem}

\begin{theorem}[Complexity of fresh-randomness bind with constant continuation]
\lean{CommunicationComplexity.PublicCoin.FiniteMessage.Protocol.rbind_complexity_const}
\label{CommunicationComplexity.PublicCoin.FiniteMessage.Protocol.rbind_complexity_const}
\leanok
\uses{CommunicationComplexity.Deterministic.FiniteMessage.Protocol.bind_complexity_const, CommunicationComplexity.Deterministic.FiniteMessage.Protocol.complexity, CommunicationComplexity.PublicCoin.FiniteMessage.Protocol, CommunicationComplexity.PublicCoin.FiniteMessage.Protocol.rbind}
\difficulty{2}
Let $p$ be a public-coin finite-message protocol over shared randomness $\Omega$,
private inputs $X$ and $Y$, and outputs $\alpha$, and let $q$ assign to each output
value $a \in \alpha$ a public-coin finite-message protocol $q(a)$ over fresh,
independent shared randomness $\Omega'$, with the same private inputs $X$ and $Y$ and
outputs in $\beta$. If there is a constant $c \in \bbn$ such that every continuation
$q(a)$ has communication complexity exactly $c$, then the protocol $\mathrm{rbind}(p,
q)$ — which runs $p$ on the $\Omega$-component of the randomness and then runs the
selected $q(a)$ on the independent $\Omega'$-component — has communication complexity
$\mathrm{complexity}(p) + c$.
\end{theorem}

\begin{definition}[Product with a deterministic protocol]
\lean{CommunicationComplexity.PublicCoin.FiniteMessage.Protocol.prodDet}
\label{CommunicationComplexity.PublicCoin.FiniteMessage.Protocol.prodDet}
\leanok
\uses{CommunicationComplexity.Deterministic.FiniteMessage.Protocol, CommunicationComplexity.Deterministic.FiniteMessage.Protocol.comap, CommunicationComplexity.Deterministic.FiniteMessage.Protocol.map, CommunicationComplexity.PublicCoin.FiniteMessage.Protocol, CommunicationComplexity.PublicCoin.FiniteMessage.Protocol.bind}
Given a public-coin protocol $p_1 : \mathrm{Protocol}\;\Omega\;X\;Y\;\alpha_1$ and a
deterministic protocol $p_2 : \mathrm{Deterministic.FiniteMessage.Protocol}\;X\;Y\;\alpha_2$,
the protocol $\mathrm{prodDet}\;p_1\;p_2$ is a $\mathrm{Protocol}\;\Omega\;X\;Y\;(\alpha_1 \times \alpha_2)$
that runs both protocols on the same inputs and pairs their outputs,
reusing the same randomness $\Omega$.
\end{definition}

\begin{theorem}[Product with a deterministic protocol runs component-wise]
\lean{CommunicationComplexity.PublicCoin.FiniteMessage.Protocol.prodDet_rrun}
\label{CommunicationComplexity.PublicCoin.FiniteMessage.Protocol.prodDet_rrun}
\leanok
\uses{CommunicationComplexity.Deterministic.FiniteMessage.Protocol, CommunicationComplexity.Deterministic.FiniteMessage.Protocol.bind_run, CommunicationComplexity.Deterministic.FiniteMessage.Protocol.comap_run, CommunicationComplexity.Deterministic.FiniteMessage.Protocol.map_run, CommunicationComplexity.Deterministic.FiniteMessage.Protocol.run, CommunicationComplexity.PublicCoin.FiniteMessage.Protocol, CommunicationComplexity.PublicCoin.FiniteMessage.Protocol.prodDet, CommunicationComplexity.PublicCoin.FiniteMessage.Protocol.rrun}
\difficulty{1}
Let $p_1$ be a public-coin finite-message protocol over shared-randomness space
$\Omega$, private input sets $X$ and $Y$, and output type $\alpha_1$, and let $p_2$ be a
deterministic finite-message protocol over $X$ and $Y$ with output type $\alpha_2$.
Write $p_1 \times p_2$ for their product: the public-coin protocol with output type
$\alpha_1 \times \alpha_2$ that runs both on the same private inputs and the same shared
randomness and pairs their outputs. Denote by $p_1(x, y; \omega)$ the output of
executing $p_1$ on private inputs $x$ and $y$ with shared randomness $\omega$, and by
$p_2(x, y)$ the output of executing the deterministic protocol $p_2$ on $x$ and $y$.
Then for every pair of private inputs $x \in X$, $y \in Y$ and every value $\omega \in
\Omega$ of the shared randomness,
\[ (p_1 \times p_2)(x, y; \omega) = \big(p_1(x, y; \omega),\ p_2(x, y)\big). \]
\end{theorem}

\begin{theorem}[Additivity of complexity under product with a deterministic protocol]
\lean{CommunicationComplexity.PublicCoin.FiniteMessage.Protocol.prodDet_complexity}
\label{CommunicationComplexity.PublicCoin.FiniteMessage.Protocol.prodDet_complexity}
\leanok
\uses{CommunicationComplexity.Deterministic.FiniteMessage.Protocol, CommunicationComplexity.Deterministic.FiniteMessage.Protocol.bind_complexity_const, CommunicationComplexity.Deterministic.FiniteMessage.Protocol.complexity, CommunicationComplexity.PublicCoin.FiniteMessage.Protocol, CommunicationComplexity.PublicCoin.FiniteMessage.Protocol.prodDet}
\difficulty{2}
Let $\Omega$ be a shared-randomness space and let $X$, $Y$ be private input sets. Let
$p_1$ be a public-coin finite-message protocol over $\Omega$, $X$, $Y$ with output type
$\alpha_1$, and let $p_2$ be a deterministic finite-message protocol over $X$, $Y$ with
output type $\alpha_2$. Form their product: the public-coin protocol over $\Omega$, $X$,
$Y$ with output type $\alpha_1 \times \alpha_2$ that runs both $p_1$ and $p_2$ on the
same inputs, sharing the randomness $\Omega$, and returns the pair of their outputs.
Then the worst-case communication complexity of this product protocol equals the sum of
the complexities of $p_1$ and $p_2$.
\end{theorem}
